\section{По време на компилация ORM архитектура}
\label{sec:architecture}

Архитектурата за обектно-релационно съпоставяне по време на компилация на библиотеката е организирана около ясно разграничение между логическия модел на данните, междинното представяне на заявките и специфичния за свързващия компонент слой за материализация. Именно този пласт носи първата основна научна ос на разработката: как компилаторът може да участва в моделирането, ограничаването и валидирането на достъпа до данни още преди фазата на изпълнение. За да бъде постигнато това, архитектурата е изградена от няколко слоя, всеки със собствена отговорност и със строго определени интерфейси.

В теоретичната глава беше показано, че унифицирането е възможно на равнището на логическия модел и междинното представяне на заявките, но не и чрез игнориране на различията между свързващите компоненти. След теорията за корутините от предходната глава тук фокусът съзнателно се връща към ядрото по време на компилация на системата: описанията на същностите, междинното представяне на заявките, \code{connector\_trait<DB>} и механизма за възможности (възможност). По този начин библиотеката се разглежда не просто като идея, а като организирана система от типове, договори и формални граници на допустимите операции.

\subsection{Архитектурни цели и ограничения}

Архитектурата е проектирана при едновременно действие на няколко ограничения. Първо, потребителският програмен интерфейс трябва да остане в рамките на идиоматичен C++, без въвеждане на отделен специфичен за домейна език (DSL) и без механизъм за отражение по време на изпълнение. Второ, структурата на заявката трябва да бъде представима като \code{constexpr} обект, така че значима част от валидирането да се извършва от компилатора. Трето, слоят на свързващия компонент трябва да бъде разширяем без наследяване от общ виртуален интерфейс, тъй като подобен интерфейс би скрил различията между компонентите и би изместил проверките към времето на изпълнение.

Четвъртото ограничение е свързано със самия характер на системата, поддържаща множество свързващи компоненти. Библиотеката трябва да може да поддържа SQL и не-SQL конектори без фалшиво уеднаквяване. Това означава, че архитектурата трябва едновременно да допуска общ програмен интерфейс и да изразява ограниченията на конкретен свързващ компонент по формален и проверим начин. Оттук естествено следва изборът на конекторен слой, базиран на характеристики, и механизъм за възможности.

\begin{figure}[H]
\centering
\begin{tikzpicture}[node distance=1.3cm and 1.4cm]
    \node[ormaccent] (app) {\textbf{Приложен слой}\\типове същности и композиция на заявки};
    \node[ormblock, below=of app] (entity) {\textbf{Слой на същностите}\\\code{property}, \code{relationship}, \code{table\_name\_trait}};
    \node[ormblock, below=of entity] (ir) {\textbf{Слой на междинното представяне}\\\code{field}, \code{Rule}, \code{select\_query}, \code{insert\_query}};
    \node[ormblock, below=of ir] (конектор) {\textbf{Конекторен слой}\\\code{connector\_trait<DB>} и labelи за възможности};
    \node[ormblock, below=of конектор] (driver) {\textbf{Слой на драйвера / време за изпълнение}\\SQL/BSON/Redis/Cypher/CQL материализация};

    \draw[ormline] (app) -- (entity);
    \draw[ormline] (entity) -- (ir);
    \draw[ormline] (ir) -- (конектор);
    \draw[ormline] (конектор) -- (driver);
\end{tikzpicture}
\caption{Слоеста архитектура на библиотеката}
\label{fig:architecture_layers}
\end{figure}

Фигура~\ref{fig:architecture_layers} показва, че библиотеката следва последователна вертикална организация. Приложният код не комуникира директно с драйвера; той работи с C++ типове и конструктори на заявки. Драйверът от своя страна не познава семантиката на домейна, а само материализираната форма, генерирана от конекторния слой. Именно поради това архитектурата може да остане едновременно строга и надграждаема.

\subsection{Слоеста организация и разделяне на отговорностите}

Най-горният слой е приложният код, който дефинира типове същности, изгражда заявки чрез \code{select}, \code{insert}, \code{update} и \code{deleteq}, и ги подава на \code{db<DB>}. Следващият слой е логическият модел на библиотеката --- описанията на същностите и междинното представяне на заявките. Третият слой е абстракцията на конектора, където \code{connector\_trait<DB>} материализира същата логика към специфичен за целевата система синтаксис и поведение. Най-долу стои конкретният драйвер, мрежов протокол (мрежов протокол) или имитираща реализация.

Тази архитектура има важен практически ефект. Ако потребителят промени \code{SQLiteDB} с \code{MongoDB}, междинното представяне на заявките остава логически същото, но конекторът го тълкува като изпълнение на филтър/проекция вместо рендиране на SQL. Това не означава, че всяко междинно представяне е допустимо върху всеки свързващ компонент; означава, че \emph{формата на абстракцията} е обща, а допустимостта се определя формално чрез механизма за възможности.

\begin{table}[H]
\centering
\small
\caption{Основни архитектурни слоеве и техните отговорности}
\label{tab:architecture_responsibilities}
\begin{tabularx}{\textwidth}{L{2.8cm}L{4.0cm}Y}
\toprule
\textbf{Слой} & \textbf{Основни артефакти} & \textbf{Отговорност} \\
\midrule
Приложен & типове същности, извиквания на конструктора на заявки & описва логиката на домейна и формулира операции в идиоматичен C++ \\
Логически & \code{property}, \code{relationship}, \code{field}, \code{Rule} & моделира данните и заявките независимо от конкретен свързващ компонент \\
Конекторен & \code{connector\_trait<DB>}, labelи за възможности & проверява допустимостта и материализира междинното представяне към поведение на целевата система \\
Време за изпълнение / драйвер & конкретни манипулатори на връзки и драйверни програмни интерфейси & изпълнява материализираната заявка и връща резултати \\
\bottomrule
\end{tabularx}
\end{table}

Таблица~\ref{tab:architecture_responsibilities} показва, че библиотеката не използва ``голям обект'', който да прави всичко. Вместо това различните решения са локализирани в отделни слоеве. Това е важно както за надграждаемостта, така и за академичния анализ, защото позволява отделните твърдения да бъдат проследени до конкретни типове и модули.

\subsection{Слоят на същностите като архитектурна основа}

Слоят на същностите е изграден върху \code{property<T, Name>}, \code{relationship<...>} и \code{table\_name\_trait<T>}. Той има двойна роля. От една страна, представя C++ домейн модела. От друга страна, служи като логическо описание на схемата или структурата на персистентните данни. Това означава, че архитектурата не прави отделен регистър по време на изпълнение от метаданни. Вместо това типовата система съдържа достатъчно информация, за да бъдат извлечени имена на полета, типове, връзки и целеви контейнер за съхранение.

Особено важен е фактът, че слоят на същностите е независим от вида база данни. Нито \code{property}, нито \code{relationship} са SQL-специфични. Те описват логически атрибути и отношения. Така се създава основата за сравнимо поведение на релационни и нерелационни свързващи компоненти. На архитектурно ниво именно това прави възможно библиотеката да твърди, че C++ моделът е първичният носител на логическата схема.

\subsection{Междинното представяне на заявките като централен архитектурен компонент}

Междинното представяне на заявките (междинно представяне на заявките) е централният архитектурен компонент. Той обединява референции към полета, предикати, заместители и допълнителни клаузи като \code{order\_by}, \code{group\_by} и \code{limit}. Вместо да се работи с текст на заявка, библиотеката изгражда типизиран обект, чийто шаблонни параметри кодират логическата структура на операцията.

От архитектурна гледна точка това решение има няколко последствия. Първо, валидирането на структурата се извършва рано. Второ, междинното представяне може да бъде анализирано от конекторния слой без загуба на семантична информация. Трето, подготвените заявки могат да запазят не SQL низ, а самото междинно представяне, което е по-естествено за библиотека, ориентирана към C++ типовата система.

Четвърто, междинното представяне на заявките е единствената зона, в която се срещат описанието на същностите и по-късната материализация в целевата система. В този смисъл то е архитектурният ``шарнир'' на цялата библиотека: достатъчно високо, за да остане независимо от съхранението, и достатъчно конкретно, за да може да бъде изпълнено.

\subsection{Конекторен слой и формалният договор}

Формално конекторният слой е реализиран чрез специализация на \code{connector\_trait<DB>}. Този подход е аналогичен на стандартни шаблони за характеристики като \code{std::iterator\_traits}. Типът \code{DB} може да бъде проста обвивка (обвивка) над манипулатор на връзка, а цялата ОРС логика за него да остане в специализацията на характеристиката. По този начин библиотеката не изисква потребителят или авторът на конектора да наследява общ базов клас.

\code{connector\_trait<DB>} определя поне три ключови аспекта: \code{wire\_type}, \code{cursor\_type} и \code{execute(...)}. По желание специализацията декларира и labelи за възможности като \code{supports\_joins}, \code{supports\_transactions}, \code{supports\_aggregation} или \code{supports\_concurrent\_execute}. Наличието или отсъствието на тези псевдо-маркери е достатъчно, за да може по-горният слой да наложи ограничения по време на компилация.

\begin{lstlisting}[language=C++,caption={Извадка от формалния договор на конектора},label={lst:connector_contract}]
template <typename DB>
концепт is_конектор = requires {
    typename конектор_trait<DB>::template wire_type<int>::type;
    typename конектор_trait<DB>::cursor_type;
};

template <typename DB, typename Cap>
inline constexpr bool has_възможност =
    detail::възможност_check<DB, Cap>::value;
\end{lstlisting}

Listing~\ref{lst:connector_contract} показва, че договорът е структурен, а не наследствен. Това е решаващ архитектурен избор. Ако съществуваше виртуален интерфейс с диспечиране по време на изпълнение (диспечиране по време на изпълнение), част от типово проверимата информация щеше да бъде изгубена. Вместо това библиотеката предпочита структура по време на компилация, при която изискванията са видими и за компилатора, и за автора на конектора.

\begin{figure}[H]
\centering
\begin{tikzpicture}[node distance=1.35cm and 1.25cm]
    \node[ormaccent] (заявка) {\textbf{Типизирана заявка}\\намерения за съединяване, транзакция, агрегация};
    \node[ormblock, below=of заявка] (db) {Фасада \code{db<DB>}};
    \node[ormblock, below left=of db] (capok) {Възможност налице\\заявката е допустима};
    \node[ormdanger, below right=of db] (capfail) {Възможност липсва\\отказ по време на компилация};
    \node[ormblock, below=of capok] (exec) {\code{connector\_trait<DB>::execute}};

    \draw[ormline] (заявка) -- (db);
    \draw[ormline] (db) -- (capok);
    \draw[ormline] (db) -- (capfail);
    \draw[ormline] (capok) -- (exec);
\end{tikzpicture}
\caption{Проверка за възможности между междинното представяне на заявката и конекторния слой}
\label{fig:capability_gating}
\end{figure}

Фигура~\ref{fig:capability_gating} е важна, защото показва къде архитектурата взема решение за допустимост. Това не е драйверът и не е логът по време на изпълнение. Решението се взема в публичния приложен програмен интерфейс (API слой), който вече има достъп до междинното представяне и до информацията за характеристиките на конектора.

\begin{table}[H]
\centering
\small
\caption{Роля на основните labelи за възможности в архитектурата}
\label{tab:capability_tags}
\begin{tabularx}{\textwidth}{L{3.6cm}Y}
\toprule
\textbf{Възможност} & \textbf{Архитектурен смисъл} \\
\midrule
\code{supports\_joins} & разрешава базираните на съединяване форми на заявки за съответния свързващ компонент \\
\code{supports\_transactions} & позволява използване на \code{transaction\_guard} и куки (куки) за \code{begin}/\code{commit}/\code{rollback} \\
\code{supports\_aggregation} & обозначава поддръжка на операции за агрегация в логическия слой на заявките \\
\code{supports\_embedding} & прави експлицитна платформено-специфичната поддръжка на вградени структури \\
\code{supports\_upsert} & маркира компоненти, при които обновяването с вмъкване (upsert) е част от естествения договор за изпълнение \\
\code{supports\_bulk\_insert} & обозначава наличието на по-ефективен път за масово записване \\
\bottomrule
\end{tabularx}
\end{table}

\subsection{Потребителският фасаден слой db<DB>}

Класът \code{db<DB>} е публичният фасаден слой. Той капсулира препратка към конкретния обект на целевата система и предоставя уеднаквен програмен интерфейс за изпълнение. На практика този слой има две роли. Първата е удобство за потребителя: библиотеката се използва чрез един и същ интерфейс независимо от свързващия компонент. Втората е архитектурен филтър: точно тук се извършват проверки по време на компилация за възможности и съвместимост.

Следователно \code{db<DB>} не е просто тънка обвивка над \code{execute}. Той е мястото, където библиотеката формализира границата между разрешено и неразрешено поведение за избрания компонент. Точно този слой позволява съобщения за грешка от вида ``конекторът не поддържа \code{JOIN} операции'', вместо грешката да се появи по-късно като неясно поведение по време на изпълнение.

\begin{lstlisting}[language=C++,caption={Извадка от фасадния слой \code{db<DB>}},label={lst:db_facade}]
template <typename DB>
    requires is_конектор<DB>
class db
{
public:
    template <typename Заявка>
    auto operator<<(Заявка q)
    {
        if constexpr (is_select_заявка<Заявка>)
        {
            if constexpr (decltype(q.join_clauses())::size > 0)
            {
                static_assert(has_възможност<DB, cap::supports_joins>,
                    "orm::db: този конектор не поддържа JOIN операции.");
            }
        }
        return конектор_trait<DB>::execute(*conn_, std::move(q));
    }
};
\end{lstlisting}

Listing~\ref{lst:db_facade} показва защо \code{db<DB>} е повече от удобна фасада. Той е формалното място, където структурата на заявката и информацията за възможностите на конектора се срещат. Именно тук архитектурата по време на компилация се превръща в реално ограничение върху употребата на програмния интерфейс.

\subsection{Миграции като архитектурно разширение}

Архитектурата включва и прототипен миграционен слой чрез \code{migrate<DB>}, \code{live\_schema}, \code{entity\_meta} и операции за дефиниране на данни (DDL) като \code{create\_table\_op}, \code{add\_column\_op}, \code{drop\_column\_op} и \code{alter\_column\_type\_op}. Този слой е особено важен от гледна точка на тезата, защото прави явна връзката между C++ модела и реалната схема на съхранение.

Миграционният слой не е универсално завършен за всеки свързващ компонент, но архитектурно показва, че библиотеката не се ограничава само до изпълнение на заявки. При достатъчно богата специализация на \code{connector\_trait<DB>} описанието на същността може да участва и в еволюцията на схемата. Това променя и самия смисъл на ОРС слоя: той вече не е само интерфейс към данните, а потенциален източник на инструментариум, съобразен със схемата (инструментариум, съобразен със схемата).

\subsection{Реализация на основните подсистеми}

След като архитектурният модел е очертан, следващият въпрос е как той се материализира в конкретни C++ конструкции и как отделните подсистеми си взаимодействат в хранилището. Настоящият раздел възстановява липсващата инженерна перспектива: не само какви слоеве съществуват, а как точно се натрупват \code{where} клаузите, как обектът на заявката се запазва за повторна употреба, как \code{migrate<DB>} превръща разликата между описанието на същността и \code{live\_schema} в DDL операции и как помощният (помощен) слой налага правила за нишкова безопасност и транзакции.

Тук се фиксира и работната нотация, без да се връща отделна речникова глава. В настоящия раздел \code{DB} означава общ тип свързващ компонент, \code{db<DB>} обозначава публичната фасада за изпълнение, \code{connector\_trait<DB>} е формалният конекторен договор, \code{field<...>} е представяне на поле по време на компилация, \code{placeholder<T>} е анонимен параметър по време на изпълнение, а \code{orm::ph<T, \_N>} е индексиран заместител. Контекстът на тази механика в хранилището е концентриран главно в \path{lib/include/ORM/entity/}, \path{lib/include/ORM/заявка/}, \path{lib/include/ORM/конектор/}, \path{lib/include/ORM/db/migration/} и във верификационния слой под \path{tests/модулни/} и \path{tests/интеграционни/}.

\subsubsection{Мини казус с типовете User и Post}

За да остане изложението непосредствено обвързано с реалния код, този раздел използва като опорен мини казус (мини казус) типовете \code{User} и \code{Post} от \path{tests/интеграционни/test_mockdb.cpp}. Те са особено полезни, защото съчетават скаларни полета, явно именуване на логическите контейнери и сценарий със съединяване (\code{JOIN}). Същият пример по-късно се използва и за свързване на заместители, поведение на подготвени заявки и утвърждавания (утвърждавания) за рендиране на SQL.

\begin{lstlisting}[language=C++,caption={Типове същности от \protect\path{tests/интеграционни/test_mockdb.cpp}},label={lst:mockdb_entities}]
struct Post
{
    orm::property<int, "id">             id;
    orm::property<int, "author_id">      author_id;
    orm::property<std::u8string, "body"> body;
};

struct User
{
    orm::property<int, "id">            id;
    orm::property<std::u8string, "name"> name;
    orm::property<double, "score">      score;
};
\end{lstlisting}

Listing~\ref{lst:mockdb_entities} показва не само как се описва моделът на същностите, но и как той непосредствено участва в тестовете. \code{table\_name\_trait<User>} и \code{table\_name\_trait<Post>} фиксират логическите контейнери \code{"users"} и \code{"posts"}, а указателите към членове (указатели към членове) като \code{\&User::id} и \code{\&Post::author\_id} по-късно участват в конструктора на заявки. Това е пряка илюстрация на един от централните мотиви на дипломната работа: моделът, формата на заявката и тестовете за изпълнение използват едни и същи артефакти по време на компилация.

\subsubsection{Скаларни свойства и именуване на контейнерите}

Реализацията на \code{property<T, Name>} свързва стойностен тип с име по време на компилация. Функцията \code{column\_name()} дава еднозначен достъп до символното име на полето и се използва от конекторите при рендиране към SQL колони, полета в документи, сегменти в ключ-стойност пространства или други специфични за целевата система понятия. Конструкцията е минималистична, но достатъчна, защото не изисква външна регистрация на метаданни по време на изпълнение.

Именуването на логическия контейнер за съхранение се реализира чрез \code{table\_name\_trait<T>}. Макар името да е наследено от ОРС традицията, архитектурно то играе по-обща роля: в SQLite, PostgreSQL и MySQL означава таблица, в MongoDB --- колекция, в Redis --- префикс на ключ, а в Neo4j --- label. Така библиотеката използва единен логически интерфейс за адресиране на физически различни структури.

Практическата последица е важна: авторът на конектора не извлича схема от външен регистър, а от типовата система. Това намалява шаблонния код и прави грешките в именуването по-лесни за локализиране както в компилационната фаза, така и в модулните (модулни) тестове.

\subsubsection{Обвивки на полета, правила и заместители}

Полето в израз на заявка се представя чрез \code{field<\&T::m>}. Това е \code{constexpr} обвивка над указател към член, която носи достатъчно информация за типа на контейнера, типа на стойността и името на колоната. Операторите върху \code{field} изграждат типизирани \code{Rule} обекти. Така изрази от вида \code{field<\&User::id> == placeholder<int>\{\}} не са низове, а описания на предикати по време на компилация.

Механизмът на заместителите (заместител) има две форми. \code{placeholder<T>} моделира анонимни параметри, които се консумират позиционно от аргументите на \code{execute()}. \code{orm::ph<T, std::placeholders::\_N>} моделира индексирани параметри, които позволяват една стойност по време на изпълнение да се използва на повече от една позиция. Това е особено важно за свързващи компоненти като SQLite, PostgreSQL и MySQL, при които слотът на заместителя е отделна част от логиката на материализация.

Тестовете \code{IndexedPlaceholderRendersQuestionMarkN}, \code{IndexedPlaceholderReusedSameArgInSQL} и \code{TwoDistinctIndexedPlaceholdersRenderCorrectSlots} в \path{tests/интеграционни/test_mockdb.cpp} показват, че тази подсистема не е теоретична добавка, а реално използван механизъм за по-точни сценарии за свързване на параметри.

\subsubsection{Изграждане на обекти на заявки}

Фабричните функции \code{select}, \code{insert}, \code{update} и \code{deleteq} създават обекти на заявки, върху които последователно се наслагват \code{where}, \code{join}, \code{order\_by}, \code{group\_by} и \code{limit}. Същественото е, че всяка от тези операции не променя низ, а връща нов типизиран обект с разширено вътрешно състояние. В резултат конструкторът на заявки остава едновременно изразителен (изразителен) на вид и статичен по природа.

Това позволява в тестовете да се проверява не само крайният резултат от рендирането, а и самата форма на междинното представяне --- например броят на \code{where} клаузите, типът на отговора или коректността на композицията на \code{GROUP BY} и \code{ORDER BY}. В \path{tests/интеграционни/test_mockdb.cpp} това се вижда в \code{SelectWithInnerJoin}, \code{SelectWithOrderByDesc}, \code{GroupByMultipleFields}, \code{SelectWithLimit} и множество други сценарии, които валидират различни фази на композицията на заявката.

Ключов имплементационен детайл е, че веригата на конструктора е съвместима с \code{constexpr}. Вътрешният помощник \code{orm::detail::tuple\_cat} използва разгъване на пакет чрез индексни последователности вместо зависимост от конкатенация по време на изпълнение. Практическата последица е, че значима част от формата на заявката може да бъде изградена и проверена още по време на компилация.

\begin{lstlisting}[language=C++,caption={Обект на заявка, деклариран като константа по време на компилация},label={lst:constexpr_query}]
using namespace orm::literals;
static constexpr auto get_active =
    orm::select(orm::field<&User::id>, orm::field<&User::name>)
        .where(orm::field<&User::score> > 0.0)
        .where(orm::field<&User::id> == orm::Заместител<int>{})
        .limit(10_per_page & 2_page);
\end{lstlisting}

Операторът \code{\&} комбинира помощниците \code{\_per\_page} и \code{\_page} в обект \code{Pagification}, чиято структура е известна по време на компилация. Хедърът \path{lib/include/ORM/заявка/limits.hpp} поддържа и симетрична форма чрез \code{*}, но използването на \code{\&} в Listing~\ref{lst:constexpr_query} е по-четливо за заявъчния стил на библиотеката. Така заявката може да бъде декларирана като \code{static constexpr} и да се използва многократно без реконструиране на формата на конструктора при всяко извикване.

\begin{figure}[H]
\centering
\begin{tikzpicture}[node distance=1.35cm and 1.2cm]
    \node[ormaccent] (api) {\textbf{Изразителен програмен интерфейс}\\\code{select}/\code{where}/\code{join}/\code{limit}};
    \node[ormblock, below=of api] (typed) {Ново типизирано\\състояние на заявката при всяка стъпка};
    \node[ormblock, below left=of typed] (shape) {Структура на заявката\\по време на компилация};
    \node[ormblock, below right=of typed] (params) {Стойности на параметрите\\по време на изпълнение};
    \node[ormblock, below=of $(shape)!0.5!(params)$] (exec) {\code{connector\_trait<DB>::execute}};

    \draw[ormline] (api) -- (typed);
    \draw[ormline] (typed) -- (shape);
    \draw[ormline] (typed) -- (params);
    \draw[ormline] (shape) -- (exec);
    \draw[ormline] (params) -- (exec);
\end{tikzpicture}
\caption{Взаимодействие между композиция на заявката, структура по време на компилация и параметри по време на изпълнение}
\label{fig:query_composition_execution}
\end{figure}

Фигура~\ref{fig:query_composition_execution} показва, че реализацията работи в два режима едновременно: структурната форма на заявката е фиксирана в типа, докато конкретните стойности се подават отделно към \code{execute(...)}. Именно това прави възможно едновременното наличие на проверки по време на компилация и параметризация по време на изпълнение.

\subsubsection{Изпълнение, резултати и достъп до полета}

На етап изпълнение \code{db<DB>} предава междинното представяне на заявката към \code{connector\_trait<DB>::execute}. Резултатът се връща като \code{orm::result<Row, FieldTuple>}. В тази конструкция \code{Row} е материализираният тип кортеж (кортеж), а \code{FieldTuple} запазва метаданните на проекцията за избраните полета. Това решение позволява както итериране върху материализираните редове, така и достъп чрез \code{get\_field<\&T::m>} или позиционен достъп чрез \code{get<I>}. Реализацията така остава типово безопасна и след изпълнението на заявката.

Тестовете \code{SelectResultRowTypeIsTuple}, \code{SelectMultiFieldRowType}, \code{SelectGetByMemberPointer} и \code{SelectGetByIndex} показват, че слоят за резултати не е просто контейнер от анонимни редове, а структурирана абстракция с типово проследима проекция. Това е важно за тезата, защото демонстрира, че типовата дисциплина не приключва с конструктора на заявки, а продължава и в обработката на резултатите (обработка на резултатите).

\subsubsection{Подготвени заявки (prepared заявки) като артефакт за многократно изпълнение}

\code{prepared\_query<DB, Query>} разширява изпълнителния модел. Вместо да кешира SQL низ, той съхранява връзка към обекта на свързващия компонент и самото междинно представяне. Това е естествено продължение на подхода по време на компилация: библиотеката третира обекта на заявката като първичен артефакт, а не като временна стъпка по пътя към текстов резултат.

\begin{lstlisting}[language=C++,caption={Извадка от \code{prepared\_query<DB, Query>}},label={lst:prepared_query}]
template <typename DB, typename Заявка>
class prepared_заявка
{
public:
    prepared_заявка(DB& conn, Заявка q) : conn_(conn), заявка_(std::move(q)) {}

    template <typename... Params>
    auto execute(Params&&... params) const
    {
        return конектор_trait<DB>::execute(
            conn_, заявка_, std::forward<Params>(params)...);
    }

    [[nodiscard]] const Заявка& заявка() const noexcept { return заявка_; }
};
\end{lstlisting}

Listing~\ref{lst:prepared_query} показва важен инженерен избор: повторната употреба е организирана около междинното представяне, а не около текстовата материализация. Това позволява една и съща логическа заявка да бъде изпълнявана многократно с различни параметри, без да се реконструира формата на конструктора. Точно това поведение се валидира от \code{PrepareReturnsExecutableObject}, \code{PrepareWithParamExecutesCorrectly}, \code{PrepareCanBeCalledMultipleTimesWithDifferentParams} и \code{PrepareExposesQueryIR}.

\begin{figure}[H]
\centering
\begin{tikzpicture}[node distance=1.35cm and 1.2cm]
    \node[ormaccent] (build) {Създаване на\\междинно представяне};
    \node[ormblock, right=of build] (prepare) {\code{db.prepare(query)}};
    \node[ormblock, right=of prepare] (store) {Съхранено междинно представяне +\\референция към \code{DB}};
    \node[ormblock, below=of store] (exec1) {\code{execute(1)}};
    \node[ormblock, left=of exec1] (exec2) {\code{execute(99)}};
    \node[ormblock, right=of exec1] (exec3) {\code{execute(7)}};

    \draw[ormline] (build) -- (prepare);
    \draw[ormline] (prepare) -- (store);
    \draw[ormline] (store) -- (exec1);
    \draw[ormline] (store) -- (exec2);
    \draw[ormline] (store) -- (exec3);
\end{tikzpicture}
\caption{Жизнен цикъл на \code{prepared\_query}: едно междинно представяне, множество изпълнения}
\label{fig:prepared_query_lifecycle}
\end{figure}

\subsubsection{Миграционен слой и DDL генериране}

Реализацията на \code{migrate<DB>} сравнява описаната от типовете същности схема с \code{live\_schema}. Когато има разлика, се създават DDL операции и се подават към \code{connector\_trait<DB>::ddl\_for(...)}. Това е добър пример за чисто разделяне на отговорностите: логиката за разлики е обща, а конкретният DDL синтаксис остава специфичен за свързващия компонент.

\begin{lstlisting}[language=C++,caption={Ядро на migration diff конвейер-а},label={lst:migration_diff}]
[[nodiscard]] std::vector<ddl_op> diff(
    const live_schema& live,
    std::span<const entity_meta> entities) const
{
    std::vector<ddl_op> ops;

    for (const auto& entity : entities)
    {
        auto live_it = live.find(entity.table_name);
        if (live_it == live.end())
        {
            ops.push_back(create_table_op{entity.table_name, entity.columns});
            continue;
        }

        for (const auto& [col, type] : entity.columns)
        {
            if (live_it->second.find(col) == live_it->second.end())
                ops.push_back(add_column_op{entity.table_name, col, type});
        }
    }
    return ops;
}
\end{lstlisting}

Тук ясно се вижда разликата между \emph{минимален работещ конектор} и \emph{пълен конектор договор}. Един свързващ компонент може да изпълнява заявки, без да поддържа DDL генериране. Но ако трябва да участва в инструментариум, съобразен със схемата, неговият договор трябва да бъде разширен с \code{ddl\_for(...)} куки и последователна миграционна верификация.

\begin{figure}[H]
\centering
\begin{tikzpicture}[node distance=1.35cm and 1.2cm]
    \node[ormaccent] (entity) {Entity metadata\\\code{entity\_meta}};
    \node[ormblock, right=of entity] (live) {По време на изпълнение schema\\\code{live\_schema}};
    \node[ormblock, below=of $(entity)!0.5!(live)$] (diff) {\code{migrate<DB>::diff}};
    \node[ormblock, below left=of diff] (ops) {DDL операции\\\code{create/add/drop/alter}};
    \node[ormblock, below right=of diff] (ddl) {\code{generate\_ddl}};
    \node[ormblock, below=of $(ops)!0.5!(ddl)$] (run) {пробно или изпълнение path};

    \draw[ormline] (entity) -- (diff);
    \draw[ormline] (live) -- (diff);
    \draw[ormline] (diff) -- (ops);
    \draw[ormline] (diff) -- (ddl);
    \draw[ormline] (ops) -- (run);
    \draw[ormline] (ddl) -- (run);
\end{tikzpicture}
\caption{Поток на schema migration подсистемата}
\label{fig:migration_pipeline}
\end{figure}

Тестовете в \path{tests/модулни/test_schema_migration.cpp} покриват създаване на таблица, добавяне и отпадане на колони, пробно кодове на завършване, генериране на DDL и поведението на автомат с крайни състояния на \code{run()}. Това е силно доказателство, че миграционният слой не е само концепция, а реално проследим конвейер.

\subsubsection{Нишкова безопасност и transaction помощници}

В библиотеката съществува отделен слой за \code{connection\_pool<DB, N>}, \code{thread\_local\_db<DB>} и \code{transaction\_guard<DB>}. Тези компоненти показват, че реализацията не спира до еднократно изпълнение на заявка, а разглежда и експлоатационни сценарии с повторна употреба на връзки, изолиране по нишки и отмяна дисциплина. Важното архитектурно решение е, че достъпът до тези механизми също е възможност-базиран. Например \code{connection\_pool} изисква \code{supports\_concurrent\_execute}, а \code{transaction\_guard} --- \code{supports\_transactions}.

\begin{lstlisting}[language=C++,caption={Извадка от нишкова безопасност помощен слоя},label={lst:нишка_safety_помощници}]
template <typename DB, std::size_t N>
class connection_pool
{
    static_assert(
        requires { typename конектор_trait<DB>::supports_concurrent_execute; },
        "connection_pool requires "
        "конектор_trait<DB>::supports_concurrent_execute. "
        "Add 'using supports_concurrent_execute = void;' "
        "to конектор_trait<DB>.");

public:
    [[nodiscard]] connection_предпазител<DB> acquire();
};

template <typename DB>
class transaction_предпазител
{
    static_assert(
        requires { typename конектор_trait<DB>::supports_transactions; },
        "transaction_предпазител requires "
        "конектор_trait<DB>::supports_transactions. "
        "Add 'using supports_transactions = void;' "
        "to конектор_trait<DB>.");
};
\end{lstlisting}

Така по време на компилация договорът не се ограничава до формата на заявката, а достига и до жизнен цикъл политиките на конектора. Това е важно за аргумента на тезата, че библиотеката не е просто конструктор за заявки, а инфраструктура за систематична и безопасна работа с свързващ компонент ресурси.

\begin{table}[H]
\centering
\small
\caption{Основни подсистеми и тяхната тестова опора}
\label{tab:core_subsystem_tests}
\begin{tabularx}{\textwidth}{L{3.1cm}L{4.1cm}Y}
\toprule
\textbf{Подсистема} & \textbf{Основни файлове} & \textbf{Тестова опора} \\
\midrule
Моделиране на същности & \path{ORM/entity/property.hpp}, \path{ORM/entity/relationship.hpp} & \path{tests/модулни/test_property.cpp}, \path{tests/модулни/test_relationship.cpp} \\
Композиция на заявки & \path{ORM/заявка/select.hpp} и свързаните заглавни файлове & \path{tests/модулни/test_заявка.cpp}, \path{tests/интеграционни/test_mockdb.cpp} \\
Подготвени заявки & \path{ORM/конектор/prepared_заявка.hpp} & Сценарии \code{Prepare*} в \path{tests/интеграционни/test_mockdb.cpp} \\
Миграция на схема & \path{ORM/db/migration/migration.hpp} & \path{tests/модулни/test_schema_migration.cpp} \\
Нишкова безопасност & \path{ORM/db/конекторs/НишкаSafety/нишка_safety.hpp} & \path{tests/модулни/test_нишка_safety.cpp} \\
\bottomrule
\end{tabularx}
\end{table}

\subsubsection{Извод за реализацията}

Основните подсистеми имат пряка опора в кода и тестовете. \path{tests/модулни/test_property.cpp} и \path{tests/модулни/test_relationship.cpp} валидират описанията на същностите. \path{tests/модулни/test_заявка.cpp} доказва, че конструкторът на заявки действително конструира правилни структури по време на компилация. \path{tests/интеграционни/test_mockdb.cpp} и \path{tests/интеграционни/test_sqlite.cpp} валидират изпълнението, сценариите за подготвени заявки и обработката на резултатите (обработка на резултатите). \path{tests/модулни/test_schema_migration.cpp} покрива логиката за разлики при миграциите, а \path{tests/модулни/test_нишка_safety.cpp} проверява слоя на жизнения цикъл за конкурентен достъп.

Следователно реализацията на основните подсистеми не е само проектно намерение, а архитектурно решение, подкрепено от систематична верификация. Това е и най-важният извод от настоящия раздел: архитектурата за обектно-релационно съпоставяне по време на компилация в хранилището не остава на равнище абстракция, а се материализира в работещи подсистеми с ясно разграничени отговорности и проследима опора в хранилището.

\subsection{Архитектурни инварианти}

На базата на разгледаните слоеве могат да се формулират няколко инварианта. Първо, приложният код никога не бива да зависи пряко от драйверен програмен интерфейс. Второ, логиката на заявката се формулира в типизирано междинно представяне, а не в низове, специфични за целевата система. Трето, разширяването към нов свързващ компонент се извършва чрез специализация на характеристика и декларации на възможности, а не чрез промяна на самия програмен интерфейс за заявки. Четвърто, ако една операция е недопустима за даден свързващ компонент, това трябва да стане видимо възможно най-рано --- за предпочитане при компилация.

Тези инварианти правят архитектурата едновременно инженерно подредена и академично защитима. Те позволяват всяко следващо разширение да бъде оценявано не само по това дали ``работи'', а и по това дали запазва вече формулираните слоеве и граници.

\subsection{Архитектурен извод}

Обобщено, архитектурата на библиотеката разделя проблема на три равнища: логически модел, типизирано междинно представяне на заявкитете и материализация в целевата система. Това разделение позволява едновременно строга типова дисциплина, формализирана разширяемост и контролирано поведение върху различни модели за съхранение. В следващата глава фокусът се премества към втората архитектурна ос --- асинхронното изпълнение --- където става ясно как същото ядро по време на компилация се свързва с жизнения цикъл на корутините, изнасянето (изнасяне) към набор от нишки и интеграцията с платформено-специфични неблокиращи интерфейси.
